% -*- coding: utf-8 -*-
\chapter{绪论}
\label{chap:intro}

\section{研究背景与科学意义}
自 1948 年克劳德·香农发表《通信的数学理论》以来，现代数字通信系统一直严格构建于经典信息论体系之上。其中，\textbf{香农信源信道分离定理（Shannon's Separation Coding Theorem）}确立了将物理通信链路划分为独立的信源编码（Source Coding，负责消除信源统计冗余）与信道编码（Channel Coding，负责按需添加校验保护冗余）的基本架构范式。在点对点平稳通信信道且假定传输码长趋向于无穷大（$n \to \infty$）的理想渐进条件下，该分离架构被严格证明能够达到信道容量极限，奠定了从 1G 到 5G 移动通信技术的物理层根基。

然而，迈入第六代移动通信系统（6G）演进周期，空天地海一体化互联、车联网（V2X）多车协同决策以及具身智能实时遥操作等极端前沿场景，对物理层传输提出了兼具“超高可靠”与“微秒级超低时延”的双重严苛约束。在这些核心场景中，传统香农分立架构遭遇了不可回避的技术死结：
\begin{enumerate}
  \item \textbf{有限码长制约与分离定理失效}：在车联网短包突发控制业务中，传输码长通常仅有数十至数百比特。此时，信源压缩与信道编码独立优化的渐进最优性彻底失效，传统级联方案的端到端率失真性能发生断崖式崩塌；
  \item \textbf{信道悬崖效应（Cliff Effect）的致命威胁}：经典数字调制体制（如 LDPC/Polar + QAM）高度依赖于严格的解调门限。当高速行驶导致瞬时信噪比骤降、低于调制编码方案（MCS）预设门限时，译码错误率急剧激增，导致整个数据包被完全丢弃，通信发生瞬间中断。
\end{enumerate}

针对这一根本性瓶颈，\textbf{语义通信（Semantic Communications）}与\textbf{深度联合信源信道编码（Deep Joint Source-Channel Coding, Deep JSCC）}应运而生。通过构建非线性可微分特征映射算子，Deep JSCC 将多媒体信息直接编码为连续信道波形，从机理上消除了量化截断与离散门限，实现了随信道状态恶化时的平滑优雅降级（Graceful Degradation），成为 6G 物理层革命性候选技术之一。

\section{国内外研究现状概述}
近年来，国内外顶尖高校与科研机构（如北京大学、帝国理工学院、清华大学、北京邮电大学等）在 Deep JSCC 领域开展了大量开拓性工作：
\begin{itemize}
  \item 在理论奠基方面，Xie 等人~\cite{xie2021deep} 率先将多头自注意力机制（Transformer）引入端到端无线收发机设计，验证了在负信噪比（$\text{SNR} \le -6\text{ dB}$）恶劣信道下其语义度量性能显著优于传统 5G 级联方案；
  \item 在连续非线性变换方面，Dai 等人~\cite{dai2022nonlinear} 提出基于变分熵模型的 NTSCC 架构，在宽信噪比动态范围内彻底消除了传统系统的悬崖效应，重构图像 PSNR 相比传统 BPG+LDPC 提升高达 3.8dB；
  \item 在工程硬件兼容性方面，Huang 等人~\cite{huang2025d2} 于 2025 年提出了离散可微分量化数字方案（$D^2$-JSCC），使得语义通信技术首次具备了在商用数字射频基带芯片上直接部署的可行性。
\end{itemize}

有关各流派的具体数学机理、横向参数对比矩阵以及目前学术界悬而未决的信噪比失配难题，将在本论文第二章中展开详尽述评。

\section{论文主要研究内容与核心创新点}
针对高速移动时变信道下信噪比动态跳变诱发的模型性能崩溃问题，本文以轻量化、自适应、硬件兼容为核心攻关导向，主要研究内容涵盖以下三个维度：
\begin{enumerate}
  \item \textbf{信噪比驱动的动态神经门控特征剪枝机制}：针对传统固定参数 Deep JSCC 无法适应信道突发跳变的缺陷，设计轻量级 CSI 门控子网络，在微秒级时间内动态剪除冗余特征通道，实现速率-算力-信噪比三维自适应平衡；
  \item \textbf{离散可微分量化与商用射频基带协同收发架构}：提出端到端梯度可反向传播的非均匀矢量量化策略，在维持平滑退化特性的前提下，实现对商用 QAM 星座图映射的硬件级兼容；
  \item \textbf{系统级仿真基准与软硬件闭环原型评测}：基于 3GPP 5G NR TDL-D 高多普勒信道模型构建完整通信链路，系统测试误码率、重构质量与硬件计算时延指标。
\end{enumerate}

\section{论文的组织结构安排}
本论文各章节的组织编排如下：
\begin{itemize}
  \item \textbf{第一章（绪论）}：阐述研究背景、科学意义、国内外研究现状及论文主要贡献；
  \item \textbf{第二章（国内外研究现状与理论基础）}：系统梳理 6G 语义通信演进脉络，剖析核心流派并给出三线对比矩阵，指出核心学术瓶颈；
  \item \textbf{第三章（面向高速时变信道的自适应 Deep JSCC 理论建模）}：详细给出所提动态门控剪枝与量化调制的数学形式化推导；
  \item \textbf{第四章（实验设置、结果分析与消融验证）}：介绍仿真平台搭建、定量对比测试与消融实验结果；
  \item \textbf{第五章（总结与未来工作展望）}：归纳全文科研成果，并指出后续在专用芯片协同设计方面的探索方向。
\end{itemize}
